% =====================================================================
% Manuscript/05_intro_story_draft.tex
%
% Agent 18 (Introduction and Historical Story) deliverable.
%
% Conventions:
%   - Canonical macros loaded from Manuscript/macros.tex (Agent 17).
%   - Citation keys are exactly the ones listed in
%     Manuscript/04_references_audit.md / Manuscript/references.bib.
%   - Notation matches Manuscript/00_notation_register.md.
%   - Theorem statements at the end of the introduction are quoted
%     verbatim (up to display formatting) from
%     Manuscript/plan1_kohler_jobin_transfer.tex Theorem~thm:KJ-FK-lower
%     /\,thm:KJ-FK-upper and
%     Manuscript/plan2_boundary_layer_assembly.tex Theorem~thm:plan2-FK.
%   - The Bernoulli spectral route (Plan~1) is mentioned only as an
%     optional remark; the centroid / space-time $\Pi$ route (Plan~2)
%     is mentioned only as an auxiliary diagnostic.
% =====================================================================

\section{Introduction}
\label{sec:intro}

% =====================================================================
\subsection{The sharp quantitative Faber--Krahn inequality}
\label{ssec:intro-headline}
% =====================================================================

Let $n\ge 2$ and let $\Om\subset\R^n$ be an open set with
$0<|\Om|<\infty$. Write $\lambda_1(\Om)$ for the first Dirichlet
eigenvalue of $-\Delta$ on $\Om$ and $\Bstar$ for the open ball
centred at the origin with $|\Bstar|=|\Om|$. The Faber--Krahn
inequality~\cite{Faber1923,Krahn1925,Krahn1926} asserts
\begin{equation}\label{eq:intro-FK}
   \lambda_1(\Om)\;\ge\;\lambda_1(\Bstar),
\end{equation}
with equality if and only if $\Om$ is a ball (up to a Lebesgue
negligible set). The scale-invariant deficit
\begin{equation}\label{eq:intro-deltaFK}
   \deltaFK(\Om)\;:=\;
   \Bigl(\tfrac{|\Om|}{\omega_n}\Bigr)^{\!2/n}\!
   \bigl(\lambda_1(\Om)-\lambda_1(\Bstar)\bigr)\;\ge\;0
\end{equation}
vanishes only at balls; the Fraenkel asymmetry
\begin{equation}\label{eq:intro-Fr}
   \Fr(\Om)\;:=\;
   \inf_{x_0\in\R^n}\frac{|\Om\,\triangle\,(\Bstar+x_0)|}{|\Bstar|}
   \;\in\;[0,2)
\end{equation}
measures the distance from~$\Om$ to a ball in $L^1$. Both proofs in
the present paper close the following sharp quantitative statement,
in the canonical normalisation fixed in
\S\ref{sec:notation}.

\begin{theorem*}[Sharp quantitative Faber--Krahn]
There is a constant $C_{\rm FK}(n)>0$, depending only on the
ambient dimension $n\ge 2$, such that for every open
$\Om\subset\R^n$ with $0<|\Om|<\infty$,
\begin{equation}\label{eq:intro-main}
   \Fr(\Om)^2\;\le\;C_{\rm FK}(n)\,\deltaFK(\Om).
\end{equation}
The exponent $2$ on $\Fr(\Om)$ is sharp, in the sense that no power
greater than $2$ on the left-hand side can hold with a constant
depending only on $n$.
\end{theorem*}

The theorem is the main result of
Brasco--De~Philippis--Velichkov~\cite{BDV} (BDV) and was reproved by
Allen--Kriventsov--Neumayer~\cite{AKN2023} via the Alt--Caffarelli--Friedman
monotonicity formula. Our contribution is to give \emph{two independent
proofs}, each with full quantitative tracking of constants, written so
that they can be read in parallel and compared step-by-step. We refer
to these as \emph{Plan~1} and \emph{Plan~2}; both proofs end at the
identical theorem statement~\eqref{eq:intro-main}, with the same
explicit dimensional constant
\begin{equation}\label{eq:intro-cFKdef}
   c_{\rm FK}(n)=1/C_{\rm FK}(n)
   \;=\;\min\!\Bigl\{
      \frac{4 L_n\,c_{\rm SV}^{\rm glob}(n)}{(n+2)\,T_n},\;
      \frac{L_n\,(2^{2/(n+2)}-1)}{4}
   \Bigr\},
\end{equation}
where $L_n=\lambda_1(B_1)$ and $T_n=T(B_1)=\omega_n/(n(n+2))$ are
dimensional (cf.\ Theorems~\ref{thm:plan1-FK-upper-intro}
and~\ref{thm:plan2-FK-intro} below).

The two proofs differ \emph{only in how the bounded sharp Saint--Venant
inequality is obtained}: Plan~1 follows the selection-principle route
of~\cite{BDV}, supplemented with detailed graph-entry, Schauder
interpolation, and nearly-spherical closure; Plan~2 follows the
level-set / metric-quotient route inspired by
Nicola--Tilli~\cite{NicolaTilli2022,GGRT2024} and uses Fusco--Julin
oscillation~\cite{FJ2014} together with a finite-perimeter first-variation
identity in a metric quotient space, an oriented radial-trace estimate,
and a boundary-layer transfer. From the global sharp Saint--Venant
inequality onwards, the two proofs share the same Kohler--Jobin
transfer~\cite{KohlerJobin1978,Brasco2014} and bounded reduction (BDV
Lemma~5.3). The aim is to make both routes explicit, with constants
tracked, so that the reader sees clearly what each route requires and
what each route delivers.

% =====================================================================
\subsection{History I: Rayleigh's conjecture, Faber, Krahn}
\label{ssec:intro-history-fk}
% =====================================================================

The inequality~\eqref{eq:intro-FK} was conjectured by Lord Rayleigh in
his \emph{Theory of Sound} (1894/96) on physical grounds: among
homogeneous membranes of given area and uniform tension, the disc
produces the lowest fundamental tone. In two dimensions Rayleigh's
conjecture was proved by Faber~\cite{Faber1923} in 1923 by
Steiner-symmetrisation of the Rayleigh quotient: rearranging the trial
function $u\in H^1_0(\Om)$ as the radially decreasing equimeasurable
$u^*$ on $\Bstar$ decreases $\int u^2$ at fixed level distribution and
decreases $\int|\nabla u|^2$ by the Pólya--Szeg\H o
principle~\cite{PolyaSzego1951}, so the Rayleigh quotient does not
increase. In 1925 Krahn~\cite{Krahn1925} extended the proof to $n=2$
(in a slightly different form) and in 1926~\cite{Krahn1926} to all
dimensions $n\ge 3$, again by symmetrisation. The argument is the
foundational application of decreasing rearrangement to spectral
geometry; modern textbook treatments are given by
Henrot~\cite{HenrotBook}, Bucur--Buttazzo~\cite{BucurButtazzo},
Bandle~\cite{Bandle1980}, Kawohl~\cite{KawohlBook} and
Kesavan~\cite{Kesavan2006}.

The rigidity statement --- equality in~\eqref{eq:intro-FK} forces $\Om$
to be a ball --- is also classical. In the modern setting of sets of
finite perimeter (introduced by De~Giorgi~\cite{DeGiorgi1958}), rigidity
in the isoperimetric inequality is the key input
(\cite[Chapter~19]{Maggi2012}); see also
Brothers--Ziemer~\cite{BrothersZiemer1988} for the equality case in the
Pólya--Szeg\H o rearrangement principle.

% =====================================================================
\subsection{History II: the Saint--Venant analogue}
\label{ssec:intro-history-sv}
% =====================================================================

A parallel functional inequality of the same vintage governs the
torsional rigidity of an elastic bar. For an open
$\Om\subset\R^n$ with $0<|\Om|<\infty$, the (variational) torsion
function $u_\Om\in H^1_0(\Om)$ is the unique minimiser of
\[
   v\;\longmapsto\;
   \tfrac{1}{2}\!\int_\Om|\nabla v|^2\,dx-\int_\Om v\,dx,
\]
equivalently $-\Delta u_\Om=1$ in $\Om$ with $u_\Om=0$ on $\partial\Om$
in the trace sense, and $u_\Om\ge 0$ a.e. The torsional rigidity is
\begin{equation}\label{eq:intro-T}
   T(\Om):=\int_\Om u_\Om\,dx
   \;=\;-2\!\int_\Om E_\Om\,dx,
\end{equation}
where, following the BDV normalisation of the notation register, the
Saint--Venant energy is
\begin{equation}\label{eq:intro-E}
   E(\Om)\;:=\;-T(\Om)/2
   \;=\;
   \min\!\Bigl\{\tfrac{1}{2}\!\int|\nabla v|^2-\!\int v\,:\,
      v\in H^1_0(\Om)\Bigr\}\le 0.
\end{equation}
Pólya and Szeg\H o~\cite{PolyaSzego1951} established the Saint--Venant
analogue of Faber--Krahn:
\begin{equation}\label{eq:intro-SV}
   T(\Om)\;\le\;T(\Bstar),
   \qquad
   E(\Om)\;\ge\;E(\Bstar),
\end{equation}
with equality only if $\Om$ is a ball (up to translations and
negligible sets). Again the proof is by decreasing rearrangement; we
adopt the BDV-normalised deficit
\begin{equation}\label{eq:intro-deltaT}
   \deltaT(\Om)\;:=\;E(\Om)-E(\Bstar)\;\ge\;0,
\end{equation}
which vanishes only at balls. Modern accounts of the Saint--Venant
inequality and its rearrangement proof are
in~\cite{KawohlBook,Bandle1980,Kesavan2006,HenrotBook,PolyaSzego1951}.
In the BDV
literature~\cite[\S2]{BDV},~\cite[\S2]{Brasco2014},~\cite[\S2]{BraDeP2017},
the quantitative form of~\eqref{eq:intro-SV} is the natural Saint--Venant
companion of Faber--Krahn, and is the form actually proved by both
of our routes before transferring to the eigenvalue side.

% =====================================================================
\subsection{The quantitative stability problem and the sharp exponent}
\label{ssec:intro-sharp-exponent}
% =====================================================================

The qualitative rigidity statement --- equality in~\eqref{eq:intro-FK}
or~\eqref{eq:intro-SV} forces $\Om$ to be a ball --- raises the
quantitative question: for $\Om$ \emph{close} to a ball, how does
$\deltaFK(\Om)$ (or $\deltaT(\Om)$) compare to $\Fr(\Om)$?

A volume-preserving ellipsoidal one-parameter family $\Om_\eps$ of
perturbations of $\Bstar$ satisfies $\Fr(\Om_\eps)\sim\eps$ while
$\lambda_1(\Om_\eps)-\lambda_1(\Bstar)\sim\eps^2$ and
$E(\Om_\eps)-E(\Bstar)\sim\eps^2$ as $\eps\to 0$; see
\cite[Introduction]{BraDeP2017} and~\cite[\S2]{BDV}. Hence the best
possible quantitative form of Faber--Krahn / Saint--Venant is
\begin{equation}\label{eq:intro-sharp-form}
   \Fr(\Om)^2\;\le\;C_{\rm FK}(n)\,\deltaFK(\Om),
   \qquad
   \Fr(\Om)^2\;\le\;C_{\rm SV}(n)\,\deltaT(\Om);
\end{equation}
the exponent $2$ is sharp. The first quantitative stability results
in the eigenvalue case~\eqref{eq:intro-FK} are due to
Hansen--Nadirashvili~\cite{HansenNadirashvili1994},
Melas~\cite{Melas1992}, and Bhattacharya~\cite{Bhattacharya2001},
all with non-sharp exponents. Fusco, Maggi, and
Pratelli~\cite{FMPLaplace}, building on the sharp quantitative
isoperimetric inequality of~\cite{FMP2008}, obtained
\begin{equation}\label{eq:intro-FMP-FK}
   \Fr(\Om)^4\;\le\;C(n)\,\deltaFK(\Om)
   \quad\text{and}\quad
   \Fr(\Om)^4\;\le\;C(n)\,\deltaT(\Om);
\end{equation}
see~\cite[Theorems~1.1, 1.4]{FMPLaplace}. The sharp
form~\eqref{eq:intro-sharp-form} is the content of~\cite{BDV} for the
ball case in any dimension $n\ge 2$; see also the survey of
Brasco--De~Philippis~\cite{BraDeP2017} and the alternative proof of
Allen--Kriventsov--Neumayer~\cite{AKN2023}. The sharp stability of
the second eigenvalue under the natural class constraint is due to
Brasco--Pratelli~\cite{BrascoPratelli2012}.

The Kohler--Jobin
inequality~\cite{KohlerJobin1978,Brasco2014} bridges the two sides of
the problem: any sharp $\deltaT$-stability translates into a sharp
$\deltaFK$-stability with the same exponent on $\Fr$. Both of our
proofs first close the sharp Saint--Venant
inequality~\eqref{eq:intro-sharp-form} and then transfer through
Kohler--Jobin. We make this bridge explicit in
Sections~\ref{ssec:intro-kj} and~\ref{ssec:intro-final-thms} below.

% =====================================================================
\subsection{The selection-principle / BDV route}
\label{ssec:intro-bdv}
% =====================================================================

The BDV selection
principle~\cite{BDV} was inspired by the prototype of
Cicalese--Leonardi~\cite{CL2012} for the isoperimetric inequality;
related second-variation selection mechanisms appear in
Acerbi--Fusco--Morini~\cite{AFM2013}. The structure of the argument
in~\cite{BDV} is: among open sets $U\subset B_R$ with $|U|=\omega_n$
and prescribed lower bound on $\Fr$, minimise a penalised functional
\[
   J_{\Om_0}(U)\;:=\;E(U)+f_{\hat\eta}(|U|)+
      \sqrt{\eps^2+\tau^2(\BDValpha(U)-\eps)^2}
\]
in which $\BDValpha$ is the BDV smooth asymmetry of
\cite[Def.~4.1]{BDV} and $f_{\hat\eta}$ is a volume penalty. The
selected minimiser $U_*$ inherits a regularity package
(\cite[Lem.~4.6, 4.9, 4.12]{BDV}: density estimates, Lipschitz
non-degeneracy, smooth Bernoulli coefficient) and via the BDV
flatness theorem (\cite[Thm~4.18]{BDV}, in the spirit of
Alt--Caffarelli~\cite{AC1981}) enters the $C^{1,\gamma}$ \emph{graph
regime}: $\partial U_*$ is the graph of a function $g$ over a sphere.
A Kinderlehrer--Nirenberg hodograph step~\cite{KN1977}, followed by
a Lieberman oblique Schauder
bootstrap~\cite{Li86a,Li86b,GT1983,BL1976}, brings $g$ into
$C^{2,\gamma_0}$ with controlled norm. The BDV nearly-spherical theorem
(\cite[Thm~3.3]{BDV}, the modern descendant of the convex-stability
estimate of Fuglede~\cite{Fuglede1989}) then closes the chain:
\[
   E(U_*)-E(B_1)\;\ge\;\tfrac{1}{32 n^2}\,\|g\|_{H^{1/2}(\partial B_1)}^2
\]
on nearly-spherical sets. After unwinding the asymmetry chain
(\cite[Lem.~4.2]{BDV}) and combining with the far-regime
$\Fr^4$-estimate (\cite[Lem.~5.1, 5.2]{BDV}, itself based on the FMP
quantitative isoperimetry of~\cite{FMP2008}), one obtains a
\emph{bounded sharp Saint--Venant} statement
\begin{equation}\label{eq:intro-SV-bounded}
   \deltaT(\Om)\;\ge\;c_{\rm SV}(n,R)\,\Fr(\Om)^2,
   \qquad
   \Om\subset B_R,\quad |\Om|=\omega_n.
\end{equation}
The dependence on the confining radius $R$ is then removed by the
constructive BDV truncation lemma~\cite[Lem.~5.3]{BDV}, applied to a
truncate $\widetilde\Om\subset B_{R_*(n)}$ whose deficit and asymmetry
are controlled by those of $\Om$. The output is a global sharp
Saint--Venant inequality $\deltaT(\Om)\ge c_{\rm
SV}^{\rm glob}(n)\Fr(\Om)^2$, which feeds into Kohler--Jobin.

This is the load-bearing Plan~1 route. The argument is conditional
only in one respect: the BDV nearly-spherical
theorem~\cite[Thm~3.3]{BDV} is invoked as a black-box quantitative
identity on the $H^{1/2}$-norm of $g$, and the regularity bootstrap
that prepares $U_*$ for that theorem requires the full BDV regularity
package. None of these inputs is conditional in the sense of being
unproved in the literature; they are simply non-trivial.

\begin{remark}[Bernoulli spectral alternative]\label{rem:intro-bernoulli}
A conceptually different conditional approach to the bounded sharp
Saint--Venant inequality~\eqref{eq:intro-SV-bounded} would invoke a
spectral / Bernoulli linearisation of the BDV selected minimiser
around the ball, using the eigenvalues of the Bernoulli operator on
$\partial B_1$. This route is sketched in some of the local source
notes (cf.\ \texttt{Plan~1/route-A-summary.md}) but is \emph{not} the
main Plan~1 proof of this manuscript: the load-bearing route is the
BDV selection principle / nearly-spherical / FMP far-regime chain
above. The Bernoulli alternative is mentioned only as a remark and
its conditional status (it requires a spectral gap statement we do
not prove here) is made explicit.
\end{remark}

% =====================================================================
\subsection{Quantitative isoperimetry background}
\label{ssec:intro-quant-iso}
% =====================================================================

The Plan~1 far-regime estimate~\cite[Lem.~5.1, 5.2]{BDV} and several
intermediate steps in Plan~2 rely on the quantitative isoperimetric
toolbox developed in the last twenty years. The cornerstone is the
sharp quantitative isoperimetric inequality of
Fusco--Maggi--Pratelli~\cite{FMP2008}: there is $C(n)>0$ such that for
every set $E\subset\R^n$ of finite perimeter with $|E|=|\Bstar|$,
\begin{equation}\label{eq:intro-FMP}
   \Fr(E)^2\;\le\;C(n)\,
   \frac{P(E)-P(\Bstar)}{P(\Bstar)},
\end{equation}
where $P$ is the De~Giorgi perimeter. The exponent $2$ is sharp. An
alternative mass-transport proof of~\eqref{eq:intro-FMP} is due to
Figalli--Maggi--Pratelli~\cite{FMP2010}; a Sobolev-type strengthening
is~\cite{FMP_BV}; an analogous sharp Sobolev / Wulff stability appears
in~\cite{CFMP2009,NeumayerWulff}; see the surveys of
Fusco~\cite{FuscoSurvey} and Maggi~\cite{Maggi2008} for context.

The strong form of Fusco--Julin~\cite{FJ2014} adds an oriented
oscillation index: for every $E\subset\R^n$ of finite perimeter with
$|E|=|\Bstar|$,
\begin{equation}\label{eq:intro-FJ}
   \Fr(E)^2+\beta(E)^2\;\le\;C(n)\,
   \frac{P(E)-P(\Bstar)}{P(\Bstar)},\qquad
   \beta(E):=\inf_{z\in\R^n}\!\Bigl(\,
      \fint_{\partial^*E}\!
      \Bigl|\nu_E(x)-\tfrac{x-z}{|x-z|}\Bigr|^2\!d\mathcal H^{n-1}
   \Bigr)^{\!1/2}.
\end{equation}
\emph{Both} the asymmetry and the normal-oscillation defect are
controlled by the isoperimetric defect, with a sharp choice of
centre. This is the load-bearing isoperimetric input for Plan~2: it
is~\eqref{eq:intro-FJ}, not~\eqref{eq:intro-FMP}, that supplies the
levelwise oscillation centres along the foliation
$\rho\mapsto E_\rho=\{u_\Om>t(\rho)\}$ with $|E_\rho|=\omega_n\rho^n$.

Two technical inputs that come with the quantitative isoperimetric
literature will be used throughout. First, the Figalli--Maggi
Sobolev--Sard theorem~\cite[App.~A, Thm~A.1]{FigalliMaggi2011}: for
the torsion function $u_\Om\in W^{2,p}_{\rm loc}$ with $-\Delta u=1$,
$\mathcal H^{n-1}(\partial^*E_t\cap\{|\nabla u|=0\})=0$ for
$\mathcal L^1$-a.e.\ level $t$; equivalently, the normal velocity
$\Vrho(x)=-t'(\rho)/|\nabla u(x)|$ is well-defined
$\mathcal H^{n-1}$-a.e.\ on $\partial^*\Erho$ for $\mathcal
L^1$-a.e.\ $\rho$. Second, the finite-perimeter / BV toolbox of
Maggi's book~\cite{Maggi2012} (coarea, divergence theorem,
slicing, structure of $\partial^*E$, rigidity in the isoperimetric
inequality) and Ambrosio--Fusco--Pallara~\cite{AFP2000} (Reshetnyak
lower semicontinuity, BV chain rule, Vol'pert slicing). The
measurable selection statement we use for the levelwise oscillation
centre on good levels is Federer~\cite[\S2.13]{Federer1969}.

% =====================================================================
\subsection{The Kohler--Jobin bridge}
\label{ssec:intro-kj}
% =====================================================================

The Faber--Krahn and Saint--Venant inequalities are not independent.
The Kohler--Jobin inequality~\cite{KohlerJobin1978} states that
\begin{equation}\label{eq:intro-KJ}
   T(\Om)\,\lambda_1(\Om)^{(n+2)/2}
   \;\ge\;
   T(\Bstar)\,\lambda_1(\Bstar)^{(n+2)/2}
\end{equation}
for every open $\Om\subset\R^n$ with $0<|\Om|<\infty$, with equality
if and only if $\Om$ is a ball. The original
proof~\cite{KohlerJobin1978} is by a rearrangement of the eigenfunction
along the level sets of the torsion function. The modern statement
used here is Brasco's~\cite[Thm~3.3]{Brasco2014}, which packages the
rearrangement as a single comparison of the functionals
$T\cdot\lambda_1^{(n+2)/2}$ across all open sets of finite
volume. Inverting the dependence at fixed volume, one obtains the
pointwise transfer
\begin{equation}\label{eq:intro-KJ-pointwise}
   \lambda_1(\Om)\;\ge\;
   \lambda_1(\Bstar)\,
   \Bigl(\tfrac{T(\Bstar)}{T(\Om)}\Bigr)^{\!2/(n+2)},
\end{equation}
and after the elementary inequality $(1-s)^{-p}-1\ge ps$ for
$s\in[0,1]$, $p\ge 0$, the Saint--Venant deficit $\deltaT(\Om)$
becomes a multiple of $\deltaFK(\Om)$ on the small-deficit regime,
with explicit dimensional multiplier $2L_n/((n+2)T_n)$
(cf.~\cite[\S3]{Brasco2014}). On the large-deficit regime
$T(\Om)\le T(\Bstar)/2$, a direct evaluation
of~\eqref{eq:intro-KJ-pointwise} provides a one-sided bound of size
$L_n(2^{2/(n+2)}-1)/4$ on $\deltaFK(\Om)/\Fr(\Om)^2$. Taking the
minimum of the two multipliers yields the explicit
constant~\eqref{eq:intro-cFKdef} for both proofs.

The Kohler--Jobin transfer is therefore the common closing step in
the manuscript: \emph{any} sharp $\deltaT$-stability gives a sharp
$\deltaFK$-stability, and the constants on the two sides are linked
by an explicit dimensional formula. Both Plan~1 and Plan~2 produce
the same input $\deltaT(\Om)\ge c_{\rm SV}^{\rm glob}(n)\Fr(\Om)^2$ to
feed into~\eqref{eq:intro-KJ}, and so the final constants
$c_{\rm FK}(n)$ on both sides
of~\eqref{eq:intro-cFKdef} agree verbatim.

% =====================================================================
\subsection{Nicola--Tilli / rearrangement / level-set viewpoint}
\label{ssec:intro-nt}
% =====================================================================

A conceptually different mechanism underlies Plan~2. The starting
point is the observation, made precise by
Nicola--Tilli~\cite{NicolaTilli2022} in the context of the
short-time Fourier transform Faber--Krahn inequality and developed
quantitatively in the stability companion of
Gómez--Guerra--Ramos--Tilli~\cite{GGRT2024}, that a Faber--Krahn-type
deficit admits an \emph{exact integral identity} over the level sets
of the relevant function. In our setting, the level sets of the
torsion function $u_\Om$ are parametrised by volume,
$\Erho:=\{u_\Om>t(\rho)\}$ with $|\Erho|=\omega_n\rho^n$,
$\rho\in[0,1]$, and the Saint--Venant deficit decomposes as
\begin{equation}\label{eq:intro-deficit-id}
   \deltaT(\Om)
   \;\simeq\;
   \int_{[\rho_*,1]}\!
   \bigl(D_H(t(\rho))+D_I(t(\rho))\bigr)\,w(\rho)\,d\rho,
\end{equation}
where $D_H(t)=\alpha(t)\beta(t)-P(\Erho)^2\ge 0$ is a
Cauchy--Schwarz defect along the level set, $D_I(t)=P(\Erho)^2-c_n
m(t)^{2-2/n}\ge 0$ is the squared-perimeter isoperimetric defect,
$\alpha(t)=\int_{\partial^*E_t}|\nabla u|^{-1}\,d\mathcal H^{n-1}$,
$\beta(t)=\int_{\partial^*E_t}|\nabla u|\,d\mathcal H^{n-1}$, and
$w(\rho)$ is an explicit non-negative weight. The identity is proved
in~\S\ref{ssec:plan2-level-set} below; it builds on the
Pólya--Szeg\H o / Talenti rearrangement
identities~\cite{PolyaSzego1951,Talenti1976,KawohlBook,Bandle1980,BBC1975,BrothersZiemer1988}
and uses the BV chain rule of Bénilan--Brezis--Crandall~\cite{BBC1975}
in the form needed by the level-set foliation.

The identity~\eqref{eq:intro-deficit-id} is fully integral: it
expresses the global deficit as an integral over the foliation of
two pointwise non-negative quantities, $D_H$ controlling the
oscillation of $|\nabla u|$ on the level set, $D_I$ controlling the
isoperimetric oscillation of the level set itself. This is the
opposite philosophy from Plan~1: instead of selecting a single
near-extremal set and analysing it as a graph over a sphere, one
analyses \emph{every} level $\rho$ of the foliation, classifies
levels as good or bad according to their isoperimetric defect, and
transfers the integrated defect bound to the metric quotient
$\X=\mathfrak M/\R^n$ in which sets are identified with their
translates.

% =====================================================================
\subsection{What Plan~1 proves and why it is included}
\label{ssec:intro-plan1}
% =====================================================================

Plan~1 follows the BDV route~\cite{BDV} step-by-step, with the level
of detail demanded by~\S2 of \texttt{MANUSCRIPT\_ASSEMBLY\_AGENTS.md}.
The skeleton is
\[
   \begin{aligned}
      &\text{selection principle (\S3.1, Agent~5)}\\
      &\quad\to\;\text{graph entry $C^{1,\gamma}$ (\S3.2, Agent~6)}\\
      &\quad\to\;\text{Schauder/interpolation $C^{2,\gamma_0}$ (\S3.2, Agent~6)}\\
      &\quad\to\;\text{BDV nearly-spherical closure (\S3.3, Agent~7)}\\
      &\quad\to\;\text{far-regime $\Fr^4$ glue (\S3.3, Agent~7)}\\
      &\quad\to\;\text{bounded sharp Saint--Venant (\S3.3, Agent~7)}\\
      &\quad\to\;\text{bounded reduction (\S3.4, Agent~8)}\\
      &\quad\to\;\text{Kohler--Jobin transfer (\S3.4, Agent~8)}\\
      &\quad\to\;\text{sharp quantitative Faber--Krahn (Theorem~\ref{thm:plan1-FK-upper-intro}).}
   \end{aligned}
\]
The bounded sharp Saint--Venant estimate
\[
   \deltaT(\Om)\;\ge\;c_{\rm SV}(n,R)\,\Fr(\Om)^2,
   \quad \Om\subset B_R,\;|\Om|=\omega_n
\]
is the substance of the BDV argument; bounded reduction is the
constructive truncation of~\cite[Lem.~5.3]{BDV}; Kohler--Jobin is
\cite[Thm~3.3]{Brasco2014}.

The reasons for including Plan~1 in this manuscript are threefold.
First, it is the canonical proof: any honest discussion of the sharp
quantitative Faber--Krahn inequality must reproduce the BDV chain,
because the structure of the rest of the literature
(\cite{BraDeP2017,AKN2023,BrascoPratelli2012,FuscoSurvey,Maggi2008})
is organised around it. Second, the constant
$c_{\rm SV}(n,R)$ produced by Plan~1 is fully explicit at every step
(selection thresholds, regularity radii, nearly-spherical constant
$c_*(n)=1/(32n^2 C_3(n))$, glue constants), and the chain shows
exactly how dimensional and confinement dependences arise. Third,
both proofs converge at the same global Saint--Venant statement and
the same Kohler--Jobin transfer; presenting Plan~1 in full makes the
common closing step explicit and makes the comparison with Plan~2
honest.

% =====================================================================
\subsection{What Plan~2 proves and why it is structurally different}
\label{ssec:intro-plan2}
% =====================================================================

Plan~2 produces \emph{the same} global sharp Saint--Venant inequality
$\deltaT(\Om)\ge c_{\rm SV}^{\rm glob}(n)\Fr(\Om)^2$ along a
structurally different route. The mechanism is a five-step chain in
the metric quotient $\X$:
\[
   \begin{aligned}
      &\text{Level-set deficit identity (\S4.1, Agent~9)}\\
      &\quad\to\;\text{Metric framework / first variation (\S4.2, Agent~10)}\\
      &\quad\to\;\text{Fusco--Julin centres and oriented radial trace (\S4.3, Agent~11)}\\
      &\quad\to\;\text{Integrated action / weighted metric trace (\S4.4, Agent~12)}\\
      &\quad\to\;\text{Boundary-layer transfer + global assembly (\S4.5, Agent~13).}
   \end{aligned}
\]
Each ingredient is necessary, and none is conditional. We summarise
the role of each.

\paragraph{(1) Exact level-set deficit identity.}
Theorem~\ref{thm:plan2-levelset-id} (Agent~9, \S\ref{ssec:plan2-level-set})
records the identity
\eqref{eq:intro-deficit-id} as a finite-perimeter equality, on the
reduced boundary $\partial^*\Erho$, with $\alpha(t)$, $\beta(t)$,
$D_H(t)$, $D_I(t)$, and $w(\rho)$ defined precisely in the notation
register. The proof uses the BV coarea theorem
(\cite[Thm~13.1]{Maggi2012},~\cite[Thm~3.40]{AFP2000}), a weak
divergence-flux identity for the gradient of the torsion function
truncated by a smooth cut-off of the radial vector field $x/|x|$, and
the BV chain rule for monotone absolutely continuous inverses
(\cite[Thm~0.4]{BBC1975}; see also~\cite[Thm~3.99]{AFP2000}). The
identity isolates two non-negative defects at every level: $D_H$ (a
Cauchy--Schwarz defect, controlling oscillation of $|\nabla u|$ on
$\partial^*E_t$) and $D_I$ (a squared-perimeter isoperimetric defect,
controlling oscillation of the level set itself).

\paragraph{(2) Metric framework and first variation.}
\S\ref{ssec:plan2-metric} (Agent~10) embeds the family
$F_\rho=\rho^{-1}\Erho$ in the quotient metric space
$(\X,\dX)=(\{\mathbf 1_A\}/\R^n, d_{L^1}/\!\sim)$ of indicator
functions modulo translation, and proves absolute continuity of
$\rho\mapsto[F_\rho]$ with metric derivative in the
Ambrosio--Gigli--Savar\'e sense~\cite[Def.~1.1.1, Thm~1.1.2]{AGS2008}.
The smooth first variation is computed under a smooth flow (with
normal velocity $\Vrho=-t'(\rho)/|\nabla u|$ ---
well-defined $\mathcal H^{n-1}$-a.e.\ on $\partial^*\Erho$ for a.e.\
$\rho$ by~\cite[Thm~A.1]{FigalliMaggi2011}), then passed to the
finite-perimeter setting through Reshetnyak lower
semicontinuity~\cite[Thm~2.38]{AFP2000}. The output is an upper bound
on the metric derivative $|\dot F_\rho|_\X$ by an explicit Lebesgue
norm of $\Vrho$ on $\partial^*\Erho$.

\paragraph{(3) Fusco--Julin centres and oriented radial trace.}
\S\ref{ssec:plan2-radial-trace} (Agent~11) chooses, on every \emph{good}
isoperimetric-defect level (i.e., where $D_I(t(\rho))$ is below a
threshold), a Borel measurable oscillation centre $z_\rho\in\R^n$
extracted from the Fusco--Julin
inequality~\eqref{eq:intro-FJ}~\cite{FJ2014}, using the
Federer measurable selection
principle~\cite[\S2.13]{Federer1969}. The asymmetry and
normal-oscillation defects of $\Erho$ at centre $z_\rho$ are then
both controlled by $D_I(t(\rho))^{1/2}$ and $D_I(t(\rho))$
respectively. The load-bearing lemma is the \emph{oriented
$L^1$-radial trace estimate}
\begin{equation}\label{eq:intro-trace}
   T_\rho
   \;=\;\!\int_{\partial^*\Erho}\!
      \Bigl|\tfrac{|x-z_\rho|}{\rho}-1\Bigr|\,d\mathcal H^{n-1}
   \;\le\;C\Bigl(
      |\Erho\triangle B_\rho(z_\rho)|
      +\!\int_{\partial^*\Erho}\!
         \Bigl|\nurho-\tfrac{x-z_\rho}{|x-z_\rho|}\Bigr|^{\!2}\!
         d\mathcal H^{n-1}
   \Bigr).
\end{equation}
The proof of~\eqref{eq:intro-trace} is a finite-perimeter divergence
identity applied to the vector field
$X(x)=\bigl|\tfrac{|x-z_\rho|}{\rho}-1\bigr|\tfrac{x-z_\rho}{|x-z_\rho|}$,
truncated near the centre to handle the $L^1_{\rm loc}$ singularity at
$z_\rho$ (Maggi~\cite[Thm~16.3]{Maggi2012}). The Fusco--Julin
inequality controls both terms on the right of~\eqref{eq:intro-trace}
by powers of $D_I$, and the level-set deficit identity converts the
integrated trace into a metric kinetic estimate.

\paragraph{(4) Integrated action and weighted metric trace.}
\S\ref{ssec:plan2-weighted-trace} (Agent~12) decomposes the foliation
into good and bad isoperimetric-defect levels, controls the bad set
by Markov's inequality applied to $D_I$, and reassembles
\[
   \int_{[\rho_*,1]}\!\bigl(\dX([F_\rho],[B_1])^2+|\dot F_\rho|_\X^2\bigr)\,\mu(d\rho)
   \;\le\;C(n,R,\rho_*)\,\deltaT(\Om),
\]
where $\mu(d\rho)=-t'(\rho)\,d\rho$. The abstract weighted metric trace
lemma (an arc-length argument in $(\X,\dX)$ with respect to the
measure $\mu$) then produces a radius $\rhad\in[\rho_*,1]$ such that
$\dX([F_{\rhad}],[B_1])^2\le C_*\,\deltaT(\Om)$ and $\rhad$ is close
to the heuristic value $\rhodelta=1-\kappa\sqrt{\deltaT(\Om)}$. The
Talenti bad-set Lebesgue estimate (\cite{Talenti1976}, used in
combination with the volume-radius parametrisation,
\cite[Thm~1.3.2]{Kesavan2006}) is what converts weighted kinetic
energy on the bad-$-t'$ set into Lebesgue kinetic energy on
$[\rho_*,1]$ without losing the $\deltaT$ rate.

\paragraph{(5) Boundary-layer transfer.}
\S\ref{ssec:plan2-boundary-transfer} (Agent~13) returns from $F_{\rhad}$ to
$\Om$. The removed boundary layer $\Om\setminus E_{\rhad}$ has
measure $O(\sqrt{\deltaT(\Om)})$, and the transfer of asymmetry from
$E_{\rhad}$ to $\Om$ squares this $\sqrt{\deltaT}$ without producing
a rate loss. The result is the bounded sharp Saint--Venant
inequality~\eqref{eq:intro-SV-bounded} in the Plan~2 form. Bounded
reduction (\cite[Lem.~5.3]{BDV}, imported through
Plan~1's \texttt{BoundedReduction}) removes the $R$-dependence, and
Kohler--Jobin~\cite{Brasco2014} closes the chain.

\begin{remark}[Auxiliary centroid / space-time route]\label{rem:intro-pi-route}
A separate centroid bound and a related space-time $\Pi$-identity
(\S\ref{ssec:plan2-aux-centroid}, Agent~14) compute, for the bulk
centroid $\zbar$, the signed radial moment
$\PiE$ and a Fubini-density identity relating
$\PiE$ across $\rho$. These quantities are auxiliary diagnostics:
the repaired Plan~2 chain above uses only the \emph{integrated}
forms of $D_H$ and $D_I$ through the
identity~\eqref{eq:intro-deficit-id} and the oriented radial trace
of~\eqref{eq:intro-trace}, not pointwise $\PiE$-control. The
$\Pi$-route is presented in the manuscript for the record and would
become load-bearing only if a separate integrated $\Pi$-control
estimate were proved (which we do not). Up to that hypothetical
upgrade, the centroid/$\Pi$-route is not used in any of the load-bearing
steps of Plan~2.
\end{remark}

The reasons for including Plan~2 are: (i) it is structurally
different from Plan~1, in that it does not select a single near-extremal
set and does not use a graph parametrisation; (ii) it exposes the
deficit as an integral over the level-set foliation, so that the role of
isoperimetric stability~\eqref{eq:intro-FJ} is direct rather than
indirect; (iii) it uses only finite-perimeter regularity throughout,
without a smoothness or graph assumption on $\partial\Om$; (iv) the
oriented $L^1$-radial trace lemma~\eqref{eq:intro-trace} is, to our
knowledge, the first place in the Faber--Krahn / Saint--Venant
quantitative literature where Fusco--Julin oscillation is used in this
form to transfer levelwise control to the metric quotient $\X$.

% =====================================================================
\subsection{Precise statement of the two final theorems}
\label{ssec:intro-final-thms}
% =====================================================================

Both proofs in this manuscript close the same theorem, in the same
normalisation, with the same explicit dimensional constant
\eqref{eq:intro-cFKdef}. We state both forms here for the record;
Plan~1 produces the proof in \S\ref{sec:plan1} and Plan~2 in
\S\ref{sec:plan2}.

\begin{theorem}[Sharp Faber--Krahn via Kohler--Jobin; Plan~1, upper form]
\label{thm:plan1-FK-upper-intro}
There is a dimensional constant $C_{\rm FK}(n)>0$ such that, for every
open set $\Om\subset\R^n$ with $0<|\Om|<\infty$,
\begin{equation}\label{eq:intro-plan1-FK}
   \Fr(\Om)^2\;\le\;C_{\rm FK}(n)\,\deltaFK(\Om),
\end{equation}
where
\[
   C_{\rm FK}(n)
   \;:=\;\frac{1}{c_{\rm FK}(n)}
   \;=\;\max\!\Bigl\{
      \frac{(n+2)\,T_n}{4 L_n\,c_{\rm SV}^{\rm glob}(n)},\;
      \frac{4}{L_n\,(2^{2/(n+2)}-1)}
   \Bigr\},
\]
with $L_n=\lambda_1(B_1)$, $T_n=T(B_1)=\omega_n/(n(n+2))$, and
$c_{\rm SV}^{\rm glob}(n)$ the global sharp Saint--Venant
constant defined in~\eqref{eq:csv-glob} of \S\ref{sec:plan1}. The
deficit $\deltaFK(\Om)$ is as in~\eqref{eq:intro-deltaFK} and
$\Fr(\Om)$ is the Fraenkel asymmetry~\eqref{eq:intro-Fr}.
Equality cases: $\deltaFK(\Om)=0$ iff $\Om$ is a ball (up to
translation and a Lebesgue negligible set); in that case
$\Fr(\Om)=0$.
\end{theorem}

\begin{theorem}[Sharp quantitative Faber--Krahn; Plan~2 final form]
\label{thm:plan2-FK-intro}
Let $n\ge 2$. There exists a constant $c_{\rm FK}(n)>0$, depending
only on $n$, given by the explicit formula~\eqref{eq:intro-cFKdef}
above, such that for every open $\Om\subset\R^n$ with $0<|\Om|<\infty$,
\begin{equation}\label{eq:intro-plan2-FK}
   \deltaFK(\Om)
   \;=\;
   \Bigl(\tfrac{|\Om|}{\omega_n}\Bigr)^{\!2/n}
   \bigl(\lambda_1(\Om)-\lambda_1(\Bstar)\bigr)
   \;\ge\;c_{\rm FK}(n)\,\Fr(\Om)^2.
\end{equation}
The exponent $2$ is sharp.
\end{theorem}

The two statements are identical up to inversion of the constant
($C_{\rm FK}(n)=1/c_{\rm FK}(n)$). The proofs end at the same place
because they share the bounded reduction~\cite[Lem.~5.3]{BDV} and the
Kohler--Jobin transfer~\cite[Thm~3.3]{Brasco2014}. The constants are
identical because both routes feed the same global Saint--Venant
constant $c_{\rm SV}^{\rm glob}(n)$ into the Kohler--Jobin transfer,
and the dimensional factors $L_n,T_n,\omega_n$ are universal.

% =====================================================================
\subsection{Related work and contributions}
\label{ssec:intro-related}
% =====================================================================

The sharp inequality~\eqref{eq:intro-main} was first
proved by Brasco--De~Philippis--Velichkov~\cite{BDV} in 2015. A
companion proof, structurally distinct in the use of the
Alt--Caffarelli--Friedman monotonicity formula, was obtained by
Allen--Kriventsov--Neumayer~\cite{AKN2023}. The sharp stability of the
second Dirichlet eigenvalue is the result of
Brasco--Pratelli~\cite{BrascoPratelli2012}. Pre-BDV non-sharp results
include Hansen--Nadirashvili~\cite{HansenNadirashvili1994},
Melas~\cite{Melas1992}, Bhattacharya~\cite{Bhattacharya2001}, and the
$\Fr^4$-bound of~\cite{FMPLaplace}. A modern survey of the quantitative
Faber--Krahn / Saint--Venant literature is in
Brasco--De~Philippis~\cite{BraDeP2017}; for the broader quantitative
isoperimetric viewpoint see Fusco~\cite{FuscoSurvey} and
Maggi~\cite{Maggi2008}.

On the isoperimetric side, the sharp quantitative
inequality~\cite{FMP2008} admits the mass-transport
proof of~\cite{FMP2010}, the BV-Sobolev strengthening of~\cite{FMP_BV},
the sharp Sobolev companion of~\cite{CFMP2009}, and the Wulff
strengthening of~\cite{NeumayerWulff}. The strong form of
Fusco--Julin~\cite{FJ2014} is what makes Plan~2 work: the oscillation
index controls a quadratic centre-dependent average of $|\nu-e_r|^2$
with the same exponent as the asymmetry. The selection-principle
mechanism that underlies Plan~1 originates in
Cicalese--Leonardi~\cite{CL2012}; the closely related second-variation
selection of~\cite{AFM2013} is in the spirit of the BDV penalised
functional but is not used directly here.

The level-set / convexity viewpoint behind Plan~2 is inspired by
Nicola--Tilli~\cite{NicolaTilli2022} and its quantitative companion
Gómez--Guerra--Ramos--Tilli~\cite{GGRT2024} in the short-time Fourier
transform setting. Our use of the metric quotient $\X=\mathfrak
M/\R^n$ and the Ambrosio--Gigli--Savar\'e metric
derivative~\cite{AGS2008} is, to our knowledge, new in the
Saint--Venant context; the oriented $L^1$-radial trace
identity~\eqref{eq:intro-trace} and its proof by truncation of the
vector field $x\mapsto||x-z|/\rho-1|(x-z)/|x-z|$ on the reduced
boundary, using~\cite[Thm~16.3]{Maggi2012}, are the technical core
that allows Fusco--Julin oscillation to be converted into integrated
metric data along the foliation.

Finally, the Kohler--Jobin
inequality~\cite{KohlerJobin1978,Brasco2014} is the
device that bridges Saint--Venant stability and Faber--Krahn
stability; it is what permits both of our proofs to share the
\emph{same} closing step and the \emph{same} final constant. We
present this bridge in full in \S\ref{ssec:plan1-final} and reuse it
verbatim in \S\ref{ssec:plan2-final}.

\paragraph{Contributions of the present paper.}
We summarise the contributions of the manuscript as follows:

\begin{enumerate}
\item A fully detailed exposition of the BDV
selection-principle / graph-entry / Schauder / nearly-spherical /
truncation chain~\cite{BDV,Brasco2014}, with explicit dimensional
constants at every step (Plan~1, \S\ref{sec:plan1}). The level of
detail follows the building blocks
\texttt{Plan~1/quantitative-selection-principle.tex},
\texttt{Plan~1/graph-entry-closure.tex},
\texttt{Plan~1/route-A-summary.tex},
\texttt{Plan~1/faber-krahn-transfer.tex}.
\item A structurally different proof of the same sharp inequality
based on an exact level-set deficit identity, a metric first-variation
in the quotient space $(\X,\dX)$ of sets modulo translation, an
oriented $L^1$-radial trace lemma anchored at Fusco--Julin oscillation
centres, a weighted metric trace, and a boundary-layer transfer. This
is Plan~2, \S\ref{sec:plan2}, and is new in the Saint--Venant /
Faber--Krahn context. The five major steps and their building blocks
\texttt{Plan~2/WIP/WIP\_LevelSetIdentity.tex},
\texttt{Plan~2/WIP/WIP\_MetricFramework.tex},
\texttt{Plan~2/WIP/WIP\_WeightedMetricTrace.tex},
\texttt{Plan~2/WIP/WIP\_BoundaryLayerTransfer.tex}, and
\texttt{Plan~2/WIP/WIP\_GlobalAssembly.tex} are integrated into the
manuscript with detail substantially greater than in the source
blocks.
\item A side-by-side bookkeeping of the constants. The bounded
Saint--Venant constants $c_{\rm SV}(n,R)$ from each route are
different intermediate quantities; the global constant
$c_{\rm SV}^{\rm glob}(n)$ that enters Kohler--Jobin is, in both
routes, given by the same expression
(\eqref{eq:intro-cFKdef}). The final constant $c_{\rm FK}(n)$
on~\eqref{eq:intro-main} is therefore identical for the two routes,
as recorded by Remark~\ref{rem:plan2-vs-plan1-normalisation} of
\S\ref{sec:plan2}.
\item A complete reference audit
(Manuscript/04\_references\_audit.md) listing every external citation
with exact bibliographic data, the theorem/proposition/page actually
used, and the manuscript location at which it enters, together with a
notation register (Manuscript/00\_notation\_register.md) fixing the
canonical symbols. No vague phrases such as `well-known' or
`standard' appear in a load-bearing step without a precise reference.
\item An internal audit pass on each of Plan~1 and Plan~2
(Manuscript/audit\_plan1.md, Manuscript/audit\_plan2.md), recording
that no conditional Bernoulli spectral input is used in the Plan~1
proof and that no centroid / space-time $\Pi$-control is used in the
load-bearing Plan~2 chain.
\end{enumerate}

\paragraph{Organisation.}
\S\ref{sec:prelims} (Preliminaries) collects the finite-perimeter,
torsion, rearrangement, and Fusco--Julin facts used in both routes.
\S\ref{sec:plan1} (Plan~1) presents the BDV selection-principle proof
from selection through to the final theorem. \S\ref{sec:plan2}
(Plan~2) presents the level-set / metric route through to the same
final theorem. The Kohler--Jobin transfer and the bounded reduction,
which are the closing steps of both routes, are written out fully in
\S\ref{ssec:plan1-final} and re-used in \S\ref{ssec:plan2-final}.
\S\ref{sec:references} (References) is the bibliography.

\bigskip

\noindent\emph{Conventions used throughout the introduction:} the
ambient dimension is $n\ge 2$, balls $B_r(z)\subset\R^n$ are open
Euclidean balls, $\Bstar$ is the open ball centred at the origin with
$|\Bstar|=|\Om|$, $\omega_n=|B_1|$, $L_n=\lambda_1(B_1)$, $T_n=T(B_1)$,
and all constants $C, c, c_{\rm SV}, c_{\rm FK}$ are positive and
depend on at most $(n,R,\rho_*)$ at every step where they appear;
dependences are declared on first use. The notation register is
\texttt{Manuscript/00\_notation\_register.md}.

% =====================================================================
% End of Manuscript/05_intro_story_draft.tex
% =====================================================================
